Đây các các blog mà tôi đã đăng trên
\href{https://www.facebook.com/groups/awsstudygroupfcj}{AWS Study
Group}.

\vspace{0.5em}\noindent\textbf{\texorpdfstring{\href{3.1-Blog1/}{Blog 1 - AWS Bedrock
Knowledge Bases: ``Mảnh ghép'' hoàn hảo cho kiến trúc Serverless
RAG}}{Blog 1 - AWS Bedrock Knowledge Bases: ``Mảnh ghép'' hoàn hảo cho kiến trúc Serverless RAG}}

Bài blog này giới thiệu về AWS Bedrock Knowledge Bases - dịch vụ giúp tự
động hóa toàn bộ quy trình RAG (chunking, embedding, indexing) mà không
cần phải tự vận hành hạ tầng Vector Database phức tạp. Việc áp dụng kiến
trúc Serverless RAG giúp hệ thống Trợ lý AI y khoa dễ dàng mở rộng và
cập nhật tri thức mới tự động thông qua Amazon S3.

\vspace{0.5em}\noindent\textbf{\texorpdfstring{\href{3.2-Blog2/}{Blog 2 - AWS Cognito:
``Mảnh ghép'' không thể thiếu cho Hệ thống Y tế
số}}{Blog 2 - AWS Cognito: ``Mảnh ghép'' không thể thiếu cho Hệ thống Y tế số}}

Bài blog này tập trung vào cách giải quyết bài toán định danh và phân
quyền đa vai trò (Bệnh nhân, Bác sĩ, Lễ tân) cho một nền tảng y tế số.
Bằng cách tích hợp AWS Cognito, dự án được tận dụng các tính năng Fully
Managed Identity, phân quyền RBAC qua JWT token và quy trình on-boarding
an toàn, giúp bảo mật dữ liệu y tế và tối ưu chi phí phát triển.

\vspace{0.5em}\noindent\textbf{\texorpdfstring{\href{3.3-Blog3/}{Blog 3 - AWS S3 Bucket:
``Mảnh ghép'' lưu trữ chuẩn Cloud-Native cho Hệ thống Y tế
số}}{Blog 3 - AWS S3 Bucket: ``Mảnh ghép'' lưu trữ chuẩn Cloud-Native cho Hệ thống Y tế số}}

Bài blog này chia sẻ kinh nghiệm chuyển đổi hệ thống lưu trữ tệp tin
(ảnh, chứng chỉ y tế, file trọng số mô hình AI) từ máy chủ cục bộ lên
Amazon S3. Việc đưa file lên S3 không chỉ giúp backend trở nên
``stateless'' và dễ dàng mở rộng, mà còn tăng cường bảo mật tuyệt đối
cho dữ liệu y tế nhạy cảm thông qua Presigned URLs.
